% CODEX REPORT::ARC-STRENGTHEN1::2026-10-08T20:24:52Z
\documentclass[11pt]{article}
\usepackage[margin=1in]{geometry}
\usepackage{amsmath,amssymb,amsthm,hyperref}
\newtheorem{theorem}{Theorem}
\newtheorem{corollary}{Corollary}
\title{ARC-STRENGTHEN1: Finite Entry, Explicit Margin, and Weaker Witness Requirements}
\author{ARC proof-strengthening agent}
\date{8 October 2026 UTC}
\begin{document}\sloppy\maketitle
\paragraph{Status.} PROVED/EXACT conditional composition; KNOWN THEOREM endpoint; OPEN SAT-specific transition. This report strengthens the alternative achievable-region route in FORMAL-SPINE1. It supplies neither a general SAT compiler nor a P versus NP separation.

\section{Exact endpoint and an explicit target}
Use the deterministic random-access model of Fortnow, Lipton, van Melkebeek and Viglas (FLMV), with encoded bit length $N$, charged index tapes and derived storage. Their Theorem 1.1 gives forbidden polynomial time-space pairs with space exponent approaching one as time exponent approaches one. More explicitly, Corollary 5.8 excludes nondeterministic linear time from
\[
 \mathrm{DTISP}[n^c,o(n^{d(c)})],\qquad
 d(c)=\frac12\left(\frac{c+2}{c^2}-c\right),\quad c>1.
\]
Theorem 3.1 transfers SAT algorithms to nondeterministic-linear-time algorithms with polylogarithmic time overhead and logarithmic extra space, computing reduction bits on demand. Thus strict exponent margins transfer this exclusion to SAT \cite{flmv}.

Define $\mathcal A$ as the set of $(t,s)$ witnessed by one uniform exact SAT machine with simultaneous bounds $N^{t+o(1)}$ and $N^{s+o(1)}$ in that model. Consequently,
\[
 t<c,\quad s<d(c)\quad\Longrightarrow\quad (t,s)\notin\mathcal A.
 \tag{1}
\]
To verify (1), choose fixed intermediate exponents between $t,c$ and between $s,d(c)$. All subpolynomial and polylogarithmic overheads fit the strict margins; the resulting simulation contradicts the cited exclusion.

In particular, the previously considered target $(6/5,1/2)$ is genuinely forbidden, despite the invalid product-line justification appearing in earlier reports. Set
\[
 c=\frac{1201}{1000}>\frac65,\qquad
 d(c)=\frac{1468676399}{2884802000}
 =\frac12+\frac{26275399}{2884802000}>\frac12.
\]
This rescues that numerical target using the actual tradeoff curve. It does not rescue any claimed construction of that target.

\section{Finite-stage closure needs no limiting algorithm}
\begin{theorem}[Existence-level finite entry]
Suppose that under SAT in P one can establish achievable pairs $(t_j,s_j)$ such that $t_j\to1$ and, for all sufficiently large $j$, $s_j\le\sigma$ for some fixed $\sigma<1$. Then $P\ne NP$.
\end{theorem}
\begin{proof}
The displayed rational function is continuous at one and $d(1)=1$. Choose $c>1$ sufficiently close to one that $d(c)>\sigma$. Choose one finite $j$ past the eventual-space threshold with $t_j<c$. Then $s_j<d(c)$ as well. The uniform machine witnessing this one pair contradicts (1).
\end{proof}
Different finite stages may have different uniform witnesses. There is no requirement of a uniformly computable stage-to-program map, a common implementation for all stages, uniform error functions across stages, or a single limiting machine. These relaxations are valid because the final proof uses just one fixed finite-stage witness. Each witness must still be correct on all inputs and have both bounds simultaneously. Programs chosen separately for individual input lengths are not permitted.

An equivalent sufficient formulation is: for one fixed $\sigma<1$, every fixed $\epsilon>0$ admits an all-input uniform machine with joint bounds $N^{1+\epsilon+o(1)}$ and $N^{\sigma+o(1)}$. Only one sufficiently small fixed $\epsilon$ is needed. This is a quantitative strengthening of the finite-entry observation already present in FORMAL-SPINE1, not a new SAT algorithm.

\section{A contraction can tolerate a nonzero error floor}
Fix an independently verified pair $(c,d(c))$ in (1). Suppose the proposed SAT-specific existence transitions give envelopes
\[
 t_{j+1}\le1+\lambda(t_j-1)+b,\qquad
 s_j\le\sigma<d(c),\qquad 0\le\lambda<1,\quad b\ge0,
 \tag{2}
\]
after a finite initial phase with finite exponent $t_0$. Every transition must actually establish membership of the next pair in $\mathcal A$.
\begin{corollary}[Tolerated local exponent loss]
If $t_*=1+b/(1-\lambda)<c$, then finitely many transitions (2) contradict (1).
\end{corollary}
\begin{proof}
Induction gives $t_j\le t_*+\lambda^j(t_0-t_*)$. If $t_0\le t_*$ the target is already reached. Otherwise the excess tends to zero, so one finite $j$ satisfies $t_j<c$. The space inequality supplies the other strict margin.
\end{proof}
Thus exact convergence to time exponent one is unnecessary. A stable floor strictly inside a verified forbidden region suffices. For $0<\lambda<1$ and $t_0>t_*$, any integer $j$ satisfying
\[
 \lambda^j(t_0-t_*)<c-t_*
\]
is sufficient. Equality is deliberately insufficient. This is an elementary envelope theorem; hypotheses (2) remain OPEN for the required general SAT domain.

\section{What known simulations and existing wrappers actually supply}
Williams's square-root-space theorem supplies a space simulation of multitape time $t$ in $O(\sqrt{t\log t})$. Its discussion explicitly identifies expensive low-degree-extension computation and asks about joint time-space tradeoffs \cite{williams}. It does not establish the joint transition (2). Even a hypothetical time-preserving map $(t,s)\mapsto(t,t/2)$ leaves the time exponent unchanged under repetition.

ARC-0008's unchanged-fallback envelope is
\[
 (t,s)\mapsto(\max(t,a),\max(s,u)).
\]
On its cofinal domain this is the identity. Cycles composed solely of these certified maps do not establish (2). This is a limitation of those envelopes, not an impossibility theorem for a better analyzed implementation. Shared setup, capped acquisition, event indexes and source-backed replay remain valid local accounting components. Their existence does not imply a fixed sublinear space exponent and contracting time envelope for general recursive search.

\section{Corpus reconciliation, tests, and handoff}
\begin{description}
\item[] CORPUS REVIEWED: local original FORMAL-SPINE1; FINAL-LINK2-ARC; ARC joint-simulation audit; repository-read ARC-0008-PREFIX-GUARD addition; repository governance snapshot and research skill. Parent lead owns current-tree and root-MAIN reconciliation. This specialized report does not claim a fresh bytewise audit of the complete repository.
\item[] DUPLICATION / CONTRADICTION AUDIT: finite witness chains, same-realization discipline, and the simulation mismatch were already present. This report adds a concrete open forbidden neighborhood for every fixed sublinear space exponent, derives the precise finite-stage quantifier relaxation, and independently validates $(6/5,1/2)$.
\item[] STALE CLAIM: earlier FINAL-LINK2-ARC and joint-simulation text used a universal $2\cos(\pi/7)$ product line. That justification remains FALSIFIED/RETIRED. No such line is used here. Timestamped originals are not modified.
\item[] WHAT CHANGED RELATIVE TO ROOT MAIN: no separation status change. Suggested main5 addition: replace an unspecified forbidden neighborhood with (1), and retain finite-stage witness freedom. No root file was edited by this delegated agent.
\item[] TESTS: companion exact-rational test verifies the endpoint margin, finite contraction formula and strict crossing, and rejects equality floors and identity-prefix progress. These are arithmetic checks, not SAT performance evidence. The source proof, not finite tests, establishes the general statements.
\item[] BARRIER / TRANSFER AUDIT: contraction potentials are an exact arithmetic transfer from dynamical systems once joint transitions are supplied. They are only an analogy for discovering those transitions. No cryptographic or random-oracle premise is used; bookkeeping supplies no nonrelativizing ingredient. No claim bypasses relativization or algebrization barriers.
\item[] NEXT TARGET: construct one all-input existence transition with a proved simultaneous space ceiling and quantitative time drift, including acquisition, failures, replay, encoding and ancestors, or find an already established SAT-specific theorem doing so. None of the primary results inspected provides that implication under SAT in P.
\end{description}
\bibliographystyle{plain}
\bibliography{2026-10-08_202452_ARC-STRENGTHEN1}
% END_OF_VERIFIED_REPORT::ARC-STRENGTHEN1
\end{document}
